\documentclass[lettersize,journal]{IEEEtran}
\usepackage{amsmath,amsfonts,amssymb}
\usepackage{algorithmic}
\usepackage{algorithm}
\usepackage{array}
\usepackage[caption=false,font=normalsize,labelfont=sf,textfont=sf]{subfig}
\usepackage{textcomp}
\usepackage{stfloats}
\usepackage{url}
\usepackage{verbatim}
\usepackage{graphicx}
\usepackage{color}
% \usepackage{appendix}
\usepackage{cite}
\newcommand{\parab}[1]{\noindent\textbf{#1} }

\hyphenation{op-tical net-works semi-conduc-tor IEEE-Xplore} 

\begin{document}
	
	\appendices
	% Insert within the manuscript's existing appendix region.
	\section{Security Proofs}
	
	\subsection{Proof of Lemma~1 (state invariant)}
	\label{app:proof-state-invariant}
	
	\emph{Proof.}
	Order the atomic transactions by their linearization points. Initially no
	response exists. Token transactions create states only while the identifier
	is fresh and return the same record on a repeated query. The first
	successful response records its context and consumes its signer's state.
	Each subsequent successful response checks equality with that recorded
	context and consumes a different live state. A mismatch or burn makes every
	remaining state unavailable, and the burned state is terminal. Every subsequent
	prescribed transaction preserves the binding
	$\mathsf{UseCtx}[\mathsf{sid}]=\mathsf{ctx}_0$, including transactions
	that burn the identifier or erase private token state.
	Induction on this transaction order proves both response assertions
	and preservation of the binding. Concurrent requests do not
	alter the argument because response release belongs to the same atomic
	transaction as consumption.\hfill$\square$
	
	\subsection{Proof of Lemma~1A (binding excludes incompatible-context acceptance)}
	\label{app:proof-context-binding}
	
	\emph{Proof.}
	An OT win requires a recorded $\mathsf{ctx}_0\ne\bot$, at least one released
	honest response for it, and a different terminal $\mathsf{ctx}^{\star}$
	under the same identifier. Lemma 1 and binding persistence imply that the
	terminal binding read returns $\mathsf{Bound}(\mathsf{ctx}_0)$, even after
	a burn. The final-state convention ensures the winning predicate and read
	refer to the same record. The state-assisted verifier accepts this read only if
	$\mathsf{ctx}^{\star}=\mathsf{ctx}_0$, a contradiction. This argument also
	covers an old active set changed to a different set and any change of the
	message or reconstructed aggregate. It does not use lattice hardness and
	makes no assertion about a stateless verifier.\hfill$\square$
	% Include in the manuscript's appendix region after its existing appendix switch.
	
	\subsection{Proof of Lemma~2 (aggregate-equivalent views)}
	\label{app:talon-lemma2-proof}
	
	{
		\emph{Proof.}
		Fix the same initial signer and identifier-service states, shared
		random-oracle functions, PRF values, call order, and external state-service
		events in both executions. Compare any two admissible mapping families
		$\mathbf M$ and $\mathbf M'$ from
		Section~\ref{subsec:flexible-active-set-selection}. In each call, both
		mappings pass the same structural and indexed-token checks and yield the
		same $C_{\mathsf{act}}$. All subsequent context and consumption checks
		therefore have identical outcomes. Whenever a response is released,
		$x_{\mathsf{ctx}}$, $c$, and
		\[
		z_j=y_j+\lambda_{\mathsf{act},j}s_jc+m_j
		\]
		depend on the mapping only through the common context. The second response
		component is the same stored $r_{j,2}$. Thus the response or rejection and
		the resulting state transition agree pointwise. Induction over the calls
		preserves the entire response vector and the state at every stopping prefix.
		For any fixed complete response collection, $\mathsf{Combine}$ uses its
		mapping only for structural validation and recomputing the common aggregate;
		its output or rejection is therefore unchanged as well. This proves exact
		joint distributional equality. The argument requires neither a common
		mapping accepted by all honest signers, equality of unrounded token sums,
		nor an assumption that different mappings are computationally difficult
		to find.\hfill$\square$
	}
	
	\subsection{Proof of Lemma~3 (mask distribution and partial responses)}
	\label{app:talon-lemma3-proof}
	
	\emph{Proof.}
	For each unordered pair of honest active signers, the difference of its two
	independent directed random-function outputs is uniform in $R_q$. The
	incidence map of the complete graph on $J$ maps these edge differences onto
	the zero-sum submodule of $R_q^{|J|}$. Surjectivity follows already from a
	spanning tree: assign its leaf edges successively to realize any prescribed
	zero-sum vector. All fibers have the same size, since this is a homomorphism
	of finite additive groups. Contributions from corrupted-pair edges translate
	the image by a vector whose sum is $-m_K$. This proves
	\eqref{eq:blue-mask-space}. Projection onto any proper set of coordinates
	is surjective with equally sized fibers. For $|J|=1$, the mask is fixed and
	there is no nonempty proper subset. Domain separation, injective context
	encoding, and Lemma 1 permit this argument independently for each new
	context; repeated requests reuse the same record.\hfill$\square$
	
	\subsection{Full Proof of Theorem~1 (identifiers without a completed context)}
	\label{app:talon-theorem1-proof}
	
	\emph{Proof.}
	Work first with ideal Gaussian sampling and the randomized honest-pair
	masks of Lemma 3. Guess an index $I\leftarrow[L]$ in the order of first
	relevant identifier-label evaluations. Call its identifier
	$\mathsf{sid}^{\dagger}$. If this index is unused, the reduction fails.
	Require the terminal identifier to equal $\mathsf{sid}^{\dagger}$ and abort
	if a response would complete a context under this identifier. Before any
	distribution change, this restriction preserves exactly
	$1/L$ of the randomized-mask game's
	incomplete-context success probability: no execution counted by
	$\mathsf{Win}_{\mathsf{PA}}^{\mathsf{core}}
	\wedge\neg\mathsf{Complete}(\mathsf{sid}^{\star})$
	ever triggers this abort when the guess is correct. This selection
	uses identifiers, including previously exposed ones, without guessing a
	message or a future active set.
	
	Every binding read, including adversarial public reads, is answered from
	the simulator's own current $\mathsf{UseCtx}$ table. The table contains
	only contexts already supplied in successful response queries; no secret
	or trapdoor is needed to answer a read. Reads neither authorize messages
	nor add labels to the $L$ count. All hybrids below maintain these tables
	and their transitions. Thus the added interface is simulated together with
	the original oracle view, including adaptive decisions based on reads.
	
	\emph{Generation gate and failure accounting.}
	We first work with the ideal capped $\mathsf{TDGen}$.
	At the first label evaluation of each nondistinguished identifier,
	run an auxiliary capped generation and abort if it returns $\bot$.
	In the initial gated hybrid, discard any successful output and
	sample the published label independently as $h\leftarrow R_q$.
	Tokens and responses are still generated by the original algorithms.
	The added abort probability and the comparison with implemented
	generation are charged once to $\epsilon_{\mathsf{badTD}}$.
	
	\emph{Screened-key hybrid.}
	Next, replace the independently uniform label by the public component
	of the successful auxiliary output $(h,\mathsf{td})$.
	Record this label consistently in the label-oracle table.
	Its successful-output law is $\mathcal D_h^{\mathsf{acc}}$.
	Retain $\mathsf{td}$ internally, but do not use it in this hybrid.
	All token and response computations still follow the original algorithms.
	
	For each label replacement, a screened-key distinguisher reproduces
	the generation gate and uses its challenge $h$ as the published label.
	It generates the signing keys itself and answers all protocol queries
	without a trapdoor for the challenge label.
	The screened-key assumption therefore bounds the computational loss
	by $\epsilon_{\mathsf{key}}$ per label, and by
	$(L-1)\epsilon_{\mathsf{key}}$ in total.
	The distinguished label remains uniform.
	
	\emph{Token and opening hybrid.}
	Keeping the generated labels and trapdoors fixed, replace each new
	nondistinguished honest token by an independent uniform
	$W_j\leftarrow R_q$.
	For comparison, the original forward sampler satisfies
	\[
	W_j=ay_j+e'_j+f_h(r_{j,1},r_{j,2}),
	\]
	where $(y_j,e'_j)\leftarrow D_{\sigma_2}^{2}$ and
	$(r_{j,1},r_{j,2})\leftarrow D_{\sigma_r}^{2}$.
	
	Replacing the image $f_h(r_{j,1},r_{j,2})$ by uniform,
	while preserving the exact conditional preimage law,
	costs at most $\delta_{\mathsf{img}}$ per token.
	Implementing that conditional sampler with $\mathsf{Pre}$
	adds at most $\delta_{\mathsf{GPV}}+d\nu$.
	Thus the joint token--opening comparison costs at most
	$\epsilon_{\mathsf{pre}}$ per token, including its later opening.
	The aggregate statistical loss is at most
	$T\epsilon_{\mathsf{pre}}$.
	
	After the image replacement, $W_j$ is uniform independently of
	$(y_j,e'_j)$.
	These variables and the opening may therefore be deferred until
	a response is needed, using
	\[
	(r_{j,1},r_{j,2})\leftarrow
	\mathsf{Pre}(h,\mathsf{td},W_j-ay_j-e'_j).
	\]
	With the ideal conditional sampler, this deferral is an exact
	distributional rewriting; its approximation error is already
	included in $\epsilon_{\mathsf{pre}}$.
	Repeated token queries return the same recorded token.
	Erased unused records need never be opened.
	
	The computational label replacement, statistical token replacement,
	and generation comparison are together charged to
	$\epsilon_{\mathsf{TD}}$.
	
	Crucially, the distinguished identifier is treated differently: sample all
	its tokens and $(y_j,e'_j,r_{j,1},r_{j,2})$ honestly under its uniform $h$.
	For each permitted response return a uniform $z_j$ and the genuinely stored
	$r_{j,2}$. Lemma 3 proves the joint distribution correct until completion,
	where this restricted game aborts. No trapdoor for the distinguished $h$
	and no equivocation of a target token is required. Terminal
	acceptance in this experiment uses $\mathsf{Vf}_{\mathsf{lat}}$;
	binding reads made by the adversary are still answered from the
	simulated service table. The same argument holds
	if its active set was chosen after seeing all its tokens.
	
	\emph{Explicit reduction to batch Hint-RLWE.}
	Let $G_{\mathsf{real}}$ denote the preceding simulation hybrid
	with $B=as+e$.
	Let $G_{\mathsf{uniform}}$ denote the same hybrid with an
	independent uniform $B$.
	Both games retain the same hidden $(s,e)$ in all hints,
	and use the same token samplers, state transitions, and abort rules.
	
	We construct a distinguisher $\mathcal B_{\mathsf{Hint}}$ whose input is
	\[
	X=\left(a,B,\{(c_i,U_i,V_i)\}_{i=1}^{Q}\right),
	\]
	drawn from either batch Hint-RLWE distribution defined above.
	The distinguisher does not receive $s$, $e$, or any honest
	Shamir share.
	It runs $\mathcal A$ and simulates its view as follows.
	
	\begin{enumerate}
		\item \emph{Initialization and corrupted keys.}
		Run $\mathcal A_0$ to obtain the static corruption set
		$\mathcal C$ and its continuation state.
		Return $0$ if the corruption set is invalid.
		Sample a fresh coefficient-wise Shamir sharing of zero,
		retaining only the shares $\{s_k\}_{k\in\mathcal C}$.
		Because $|\mathcal C|<t$, their joint distribution is identical
		to the restriction of a sharing of any secret $s$.
		
		Generate the PRF material and randomized honest--honest masks
		as in the preceding hybrid.
		Give $\mathcal A_1$ the corrupted secret-key material and
		\[
		\mathsf{vk}=(a,\lfloor B\rceil_{p_1},t,n,a_{\max}).
		\]
		Follow the same distinguished-label selection convention
		and initialize all oracle and identifier-state tables.
		
		\item \emph{Labels and challenge queries.}
		Sample the distinguished label uniformly.
		For every nondistinguished label, run the capped
		$\mathsf{TDGen}$ and retain its successful pair
		$(h,\mathsf{td})$.
		Return $0$ on a generation abort.
		Repeated label queries receive the same recorded output.
		
		Before the terminal output, assign the $i$th input hint
		to the $i$th distinct challenge-oracle input,
		and answer that input with $c_i$.
		Repeated inputs reuse their assigned index.
		Index hints by oracle inputs, rather than by challenge values.
		Ordinary queries to the remaining random-oracle domain
		are handled by lazy sampling.
		
		If terminal verification makes a fresh challenge query,
		answer it with an independent uniform challenge.
		No hint is required for this terminal-only query.
		
		\item \emph{Token queries.}
		Apply the token-state checks and idempotence rules
		of the preceding hybrid.
		For a new honest token under the distinguished identifier,
		sample
		\[
		(y_j,e'_j)\leftarrow D_{\sigma_2}^{2},
		\quad
		(r_{j,1},r_{j,2})\leftarrow D_{\sigma_r}^{2},
		\]
		and compute
		\[
		W_j=ay_j+e'_j+f_h(r_{j,1},r_{j,2}).
		\]
		Store the corresponding private record.
		
		For a new honest token under a nondistinguished identifier,
		sample $W_j\leftarrow R_q$ and store its unopened record.
		In both cases, return the public token
		$(j,\mathsf{sid},W_j)$.
		
		\item \emph{Response queries.}
		Perform the structural, recorded-token, aggregate,
		challenge, and bound-context checks of the preceding hybrid.
		Maintain the same rejection and burn behavior.
		
		For the current context, use $J$, $s_K$, and $m_K$
		as defined in Lemma~3.
		The values $s_K$ and $m_K$ are computable from the
		corrupted key material.
		
		Under the distinguished identifier, return
		\[
		\sigma_j=(z_j,r_{j,2}),
		\quad z_j\leftarrow R_q,
		\]
		using the genuinely stored $r_{j,2}$.
		Return $0$ if a successful response would complete
		the distinguished context.
		
		For a nondistinguished nonfinal honest responder,
		sample $(y_j,e'_j)\leftarrow D_{\sigma_2}^{2}$,
		choose $z_j\leftarrow R_q$, and sample
		\[
		(r_{j,1},r_{j,2})\leftarrow
		\mathsf{Pre}(h,\mathsf{td},W_j-ay_j-e'_j).
		\]
		Store $y_j$ and $z_j$, and return $(z_j,r_{j,2})$.
		
		For the last honest responder $j$, retrieve the hint
		assigned to the context's challenge input and compute
		\begin{equation}
			\begin{aligned}
				w_j^{\mathsf{res}}
				&=W_j-aU_i+Bc_i-V_i,\\
				(r_{j,1},r_{j,2})
				&\leftarrow
				\mathsf{Pre}(h,\mathsf{td},w_j^{\mathsf{res}}),\\
				z_j
				&=U_i-s_Kc_i-m_K
				+\sum_{k\in J\setminus\{j\}}(y_k-z_k).
			\end{aligned}
			\label{eq:app-talon-hint-simulation}
		\end{equation}
		Return $(z_j,r_{j,2})$.
		Each successful response performs the prescribed atomic
		binding, consumption, erasure, and $\mathsf{RespSet}$ update.
		
		\item \emph{Binding and burn queries.}
		Answer binding reads from the current $\mathsf{UseCtx}$ table.
		Process burn queries using the same erasure and persistence rules.
		These operations require no honest secret share.
		
		\item \emph{Terminal decision.}
		On $\mathcal A$'s terminal output, evaluate
		$\mathsf{Win}_{\mathsf{PA}}^{\mathsf{core}}$
		using $\mathsf{Vf}_{\mathsf{lat}}$ and the recorded state.
		Output $1$ exactly when no abort occurred,
		the terminal identifier is distinguished, and
		\[
		\mathsf{Win}_{\mathsf{PA}}^{\mathsf{core}}
		\wedge
		\neg\mathsf{Complete}(\mathsf{sid}^{\star})
		\]
		holds.
		Otherwise output $0$.
	\end{enumerate}
	
	To verify the simulation for a real-distribution input,
	write $U_i=y_i+sc_i$ and $V_i=e'_i+ec_i$.
	When $B=as+e$, the residual satisfies
	\[
	\begin{aligned}
		W_j-aU_i+Bc_i-&V_i
		=W_j-a(y_i+sc_i)+(as+e)c_i\\
		&-(e'_i+ec_i) =W_j-ay_i-e'_i.
	\end{aligned}
	\]
	Thus the deferred Gaussian pair of the last responder
	can be identified with $(y_i,e'_i)$.
	Its response satisfies
	\[
	z_j
	=y_i+(s-s_K)c_i-m_K
	-\sum_{k\in J\setminus\{j\}}z_k
	+\sum_{k\in J\setminus\{j\}}y_k,
	\]
	which is exactly \eqref{eq:blue-last-response}.
	
	The uniform token revealed no information about this deferred pair,
	and the signer had not previously released a response.
	Hence this identification remains valid when the final responder
	is chosen adaptively.
	Since an identifier has only one activated context,
	one challenge input closes at most one honest-response vector.
	
	Consequently, a real-distribution input reproduces
	$G_{\mathsf{real}}$, whereas a uniform-distribution input
	reproduces $G_{\mathsf{uniform}}$.
	Let $p_{\mathsf{real}}$ and $p_{\mathsf{uniform}}$
	be their respective probabilities of the terminal decision event.
	By construction,
	\[
	\begin{aligned}
		|p_{\mathsf{real}}-p_{\mathsf{uniform}}|
		=
		\operatorname{Adv}^{\mathsf{Hint\text{-}RLWE}}
		_{\mathcal B_{\mathsf{Hint}}} \leq \epsilon_{\mathsf{Hint}}.
	\end{aligned}
	\]
	The distinguisher has polynomial overhead under the stated
	generator and sampler contracts.
	The earlier label and token simulation errors remain
	accounted for separately in $\epsilon_{\mathsf{TD}}$.
	
	It remains to bound successful outputs in this uniform-key game. Embed a
	ring-SIS instance $(a,B,h)$ by using its first two coordinates in the
	verification key and its third as the distinguished label. The simulator
	for this game can generate its own hidden dummy $(s,e)$ and all hints,
	independently of $B$, and is therefore executable and rewindable without
	an external hint oracle. Other labels use the stored trapdoors.
	
	For the forking argument, regard the adversary and its simulator
	as one randomized program $\mathcal F$.
	Separate its non-challenge randomness $\rho$ from the challenge
	answers $c_1,\ldots,c_Q$.
	The tape $\rho$ includes the adversary's coins, the label-selection
	coins, all token and sampler coins, and the locally generated
	dummy secret $(s_0,e_0)$ and Gaussian hint pairs $(y_i,e'_i)$.
	
	For each distinct challenge input indexed by $i$, compute
	\[
	U_i=y_i+s_0c_i,\qquad V_i=e'_i+e_0c_i.
	\]
	Thus changing a challenge answer also recomputes its associated hint,
	using the same dummy secret and Gaussian base pair.
	
	On a successful restricted forgery with a previously evaluated
	terminal challenge input, $\mathcal F$ returns its query index
	and the terminal message--signature pair.
	Otherwise it returns index zero.
	
	If the terminal challenge input was never evaluated before the output,
	its answer, sampled during verification, equals the proposed $c$ with
	probability at most $1/M$. Otherwise, let $k^\star\in[Q]$ be the index of the first
	evaluation of the terminal challenge input.
	During the first execution, record a snapshot before answering
	each new challenge input.
	The snapshot for $k^\star$ is taken after the input and its
	associated label have been determined, but before the challenge
	answer is recorded or returned.
	
	This snapshot contains the adversary and simulator configurations,
	including any pending oracle invocation and its local variables.
	It also contains both random-oracle tables, the materialized hint
	records, all identifier-service tables, token and private-state
	records, trapdoor records, counters, and random-tape positions.
	In particular, it preserves all consumption and erasure flags.
	If the challenge evaluation occurs inside a response query,
	the snapshot includes that suspended response computation.
	
	Run the second execution with the same non-challenge tape $\rho$
	and the same answers $c_1,\ldots,c_{k^\star-1}$.
	Independently sample new answers
	$c'_{k^\star},\ldots,c'_Q$.
	This can be implemented by restarting $\mathcal F$ from its initial
	state and replaying the common prefix.
	Equivalently, restore the snapshot and resume the pending challenge
	evaluation with $c'_{k^\star}$.
	
	Only the prefix tables are restored.
	All records created after the snapshot in the first execution
	are discarded before the second continuation.
	Each continuation maintains its own oracle and service state,
	and recomputes its hints from its current challenge answers.
	The no-completion abort rule remains in force in both executions.
	
	Let $k'$ be the terminal challenge index returned by the second
	execution.
	Declare the fork successful only if both executions yield successful
	restricted forgeries, $k'=k^\star$, and
	$c'_{k^\star}\ne c_{k^\star}$.
	Otherwise the forking algorithm fails.
	
	Because the executions agree before the $k^\star$th answer,
	the corresponding query input is identical in both executions.
	Canonical injective encoding therefore fixes
	$(\mathsf{vk},\mathsf{sid}^{\dagger},
	\mathsf{act}^{\star},C^{\star},\mathsf{msg}^{\star})$.
	The associated label $h$ is already fixed by the internal
	label-evaluation convention.
	A successful fork consequently yields two norm-valid signatures
	for the same complete challenge input, with different challenges.
	Lemma~4 then extracts a nonzero ring-SIS solution.
	
	All restored state belongs to the reduction's simulated experiment.
	No rollback interface for the deployed binding service is assumed.
	
	For completeness, let $p$ be the success probability in the uniform-key
	game restricted to outputs with an already evaluated challenge input.
	The classical forking inequality
	gives a solver with success at least
	\[
	p\left(\frac pQ-\frac1M\right).
	\]
	Hence $p\leq Q/M+\sqrt{Q\epsilon_{\mathsf{SIS}}}$. Adding the unqueried
	case, undoing the two hybrid changes and the $1/L$ selection, and charging
	the initial PRF and whole-run sampling errors proves
	\eqref{eq:blue-incomplete-theorem}. All losses count the full adversarial
	interaction; no assumption on a fixed subset or a fixed target message was
	introduced.\hfill$\square$
	
	\subsection{Proof of Lemma~4 (fixed-input extraction with scaled lifting)}
	\label{app:talon-lemma4-proof}
	
	\emph{Proof.}
	In $\mathcal M_q$, define
	\[
	\begin{aligned}
		X&=az+hr-\mathsf{Lift}_1(bc_{\mathbb Z}),\\
		X'&=az'+hr'-\mathsf{Lift}_1(bc'_{\mathbb Z}),\\
		\eta_B&=B-\mathsf{Lift}_1(b).
	\end{aligned}
	\]
	The lift is an additive homomorphism compatible with multiplication by
	integer polynomials modulo $q\mathbb Z^N$. Thus no equality between
	elements of different finite rings is assumed. Since
	$\|\eta_B\|_\infty\leq q/(2q_1)$ and both reconstructions equal $C$,
	\[
	\lfloor X\rceil_{p_2}-\lfloor X'\rceil_{p_2}
	=\Delta'-\Delta.
	\]
	Applying the rounding error bound to $X$ and $X'$ separately gives
	\[
	\|X-X'\|_2
	\leq\frac q{q_2}\left(2\beta_\Delta+\sqrt N\right).
	\]
	All these norms are quotient norms, evaluated using shortest centered
	representatives; the triangle inequality is valid in the quotient.
	Let $u=z-z'$, $d=c-c'$, and $v=r-r'$. Then
	\[
	au-Bd+hv=X-X'-\eta_Bd\quad\text{in }\mathcal M_q.
	\]
	The left side lies in $R_q$. Since $\|d\|_1\leq2\tau$, Young's
	inequality gives
	\[
	\|\eta_Bd\|_2\leq\frac{q\tau}{q_1}\sqrt N.
	\]
	Take $\delta$ to be the centered representative of $-(au-Bd+hv)$.
	Then
	\[
	[\,1\ \ a\ \ B\ \ h\,]
	\begin{pmatrix}\delta\\u\\-d\\v\end{pmatrix}=0,
	\qquad \|\delta\|_2\leq\beta_\delta.
	\]
	Moreover, $\|u\|_2\leq2\beta_z$, $\|v\|_2\leq2\beta_r$, and
	$\|d\|_2\leq2\sqrt\tau$. The vector is nonzero because $c\ne c'$
	as canonically encoded ring elements. Summing these squared bounds proves
	\eqref{eq:blue-extraction-bound}.\hfill$\square$
	
	\subsection{Full Proof of Theorem~2 (state-assisted preprocessing-aware security)}
	\label{app:talon-theorem2-proof}
	
	\emph{Proof.}
	The generic lifting theorem bounds the
	state-assisted PA advantage by the lattice-only incomplete-context
	advantage. Theorem~1 bounds that advantage by the displayed
	right-hand side. The target-completion abort in Theorem~1 discards no
	execution counted by its incomplete-context experiment; the generic
	binding argument handles completed target identifiers.
	\hfill$\square$
	% Append after the existing Appendix C in the manuscript's appendix region.
	% Do not add another \appendix or \appendices command.
	
	\subsection{Screened Generation and Its Output Law}
	\label{par:blue-screened-generation}
	Put $R=\mathbb Z[X]/(X^N+1)$, $d=2N$, and $\gamma=1.17\sqrt q$.
	A trial samples independent integer coefficient vectors
	$f,g\leftarrow D_{\sigma_{\mathsf{td}}}$ before reduction modulo $q$.
	Reject unless $f\bmod q$ is invertible. Run a fixed implementation of
	Falcon's NTRU equation solver on $(f,g,q)$, with a fixed polynomial
	operation and output-bit-length budget; a timeout or failure rejects
	the trial.\footnote{Falcon specification v1.2, Algorithm~6;
		\url{https://falcon-sign.info/falcon.pdf}. The additional exact
		certification and resource limits below are part of our specified
		screening rule, not a claimed benchmark of Falcon at these parameters.}
	For a returned pair $(F,G)$, verify exactly in $R$ that $fG-gF=q$.
	Set $h=gf^{-1}\bmod q$, and form the integer coefficient basis
	$\mathcal B$ from the negacyclic rotations of $(g,-f)$ and $(G,-F)$.
	These vectors generate
	\[
	\Lambda_h=\{(u,v)\in\mathbb Z^{2N}:u+hv=0\bmod q\}.
	\]
	Indeed their determinant has absolute value $q^N$, the index of
	$\Lambda_h$, and each vector lies in this kernel. Compute the squared
	Gram--Schmidt lengths of $\mathcal B$ by exact rational arithmetic,
	or use certified interval upper bounds. Accept only if
	\[
	\max_{i\in[d]}\|\widetilde b_i\|_2^2
	\leq \gamma^2=\frac{13689}{10000}q.
	\]
	An inconclusive certificate rejects. Thus every accepted basis satisfies
	the claimed GS bound by construction, without an assumption about the
	typical norm of sampled $(f,g)$. This is a specification of a certified
	generator; no implementation or acceptance-rate measurement at
	TT-128-R1 is claimed here.
	
	Let $\mathsf{Acc}$ include every test and resource limit above, and let
	$p_{\mathsf{acc}}=\Pr[\mathsf{Acc}]$ for an independent trial, including
	any solver randomness. With a trial cap $J_{\mathsf{gen}}$, return the
	first accepted $(h,\mathcal B)$ or $\bot$ if all trials reject. Then
	\[
	\begin{aligned}
		\Pr[\mathsf{TDGen}\ne\bot]
		&=1-(1-p_{\mathsf{acc}})^{J_{\mathsf{gen}}},\\
		\mathbb E[\#\text{trials}]
		&=\frac{1-(1-p_{\mathsf{acc}})^{J_{\mathsf{gen}}}}
		{p_{\mathsf{acc}}}.
	\end{aligned}
	\]
	Without the cap the expectation is $1/p_{\mathsf{acc}}$; if one trial
	has expected cost at most $C_{\mathsf{gen}}$, the expected total cost
	is at most $C_{\mathsf{gen}}/p_{\mathsf{acc}}$.
	\emph{Generation assumption.} For the fixed solver, budgets and
	screening rule, assume a stated inverse-polynomial lower bound
	$p_{\mathsf{acc}}\geq p_0>0$ and polynomial trial cost. No numerical
	$p_0$ has been established for TT-128-R1. Bounded-time Gaussian
	implementations are compared with this ideal sampling experiment using
	the per-generator error $\delta_{\mathsf{gen}}$ below.
	A concrete instantiation must fix the solver version, arithmetic
	precision, resource limits, trial cap and a justified $p_0$ before
	claiming the resulting generation cost. The formulas above do not
	substitute for these missing concrete choices.
	
	Conditioned on success, the output distribution is exactly the pushforward
	of the trial distribution conditioned on $\mathsf{Acc}$, independently
	of the trial cap. In particular, define the public-key law
	\[
	\mathcal D_h^{\mathsf{acc}}
	=\mathcal L(gf^{-1}\bmod q\mid\mathsf{Acc}).
	\]
	\emph{Screened-key assumption.} For every distinguisher of the time and
	query budget used by the reduction, assume
	\[
	\left|\Pr[\mathcal D(h)=1:h\leftarrow\mathcal D_h^{\mathsf{acc}}]
	-\Pr[\mathcal D(h)=1:h\leftarrow R_q]\right|
	\leq\epsilon_{\mathsf{key}}.
	\]
	The trapdoor and trial history are not given to this distinguisher.
	This is an explicit assumption on the \emph{screened} distribution,
	not a consequence proved here from unscreened NTRU key indistinguishability.
	Filtering depends on secret coefficients and solver behavior, so even a
	lower bound on its acceptance probability does not by itself justify
	transferring pseudorandomness through that filter.
	
	\subsection{Preimages for Every Target and Statistical Bounds}
	\label{par:blue-preimage-contract}
	Use the GPV randomized nearest-plane sampler as a reference
	implementation of $\mathsf{Pre}$~\cite{GPV08}.
	For $w\in R_q$, set $t_w=(\overline w,0)\in\mathbb Z^{2N}$.
	Starting with center $t_w$, for $i=d,\ldots,1$ compute the current
	center's projection coefficient onto $\widetilde b_i$, sample an integer
	with that center and width $\sigma_r/\|\widetilde b_i\|_2$, and subtract
	its multiple of $b_i$ from the current center. If $v\in\Lambda_h$ is
	the sum of those multiples, output $r=t_w-v$, reduced modulo $q$ only
	at the protocol interface. Then $f_h(r)=w$ exactly.
	
	For any accepted basis and any $w$, the target conditional distribution is
	$D_{t_w+\Lambda_h,\sigma_r}$ on integer vectors. The GPV reference sampler approximates this distribution. Assume a one-dimensional
	sampler whose distance from its exact shifted discrete Gaussian is at
	most $\nu$ uniformly over every center and width encountered; any
	finite-precision or truncation errors must be included in this bound.
	Let $0<\epsilon<1$ and require
	\[
	\sigma_r\geq\gamma
	\sqrt{\frac{\ln(2d(1+1/\epsilon))}{\pi}}.
	\]
	The GS smoothing bound gives
	$\sigma_r\geq\eta_\epsilon(\Lambda_h)$ and every one-dimensional
	width is at least $\eta_\epsilon(\mathbb Z)$. Gaussian coset-mass
	bounds and the sequential sampler's normalization factors then give
	the conservative bounds
	\[
	\begin{aligned}
		\delta_{\mathsf{img}}=\frac{2\epsilon}{1-\epsilon}&,\,
		\delta_{\mathsf{GPV}}=\left(\frac{1+\epsilon}{1-\epsilon}\right)^d-1,\\
		\sup\nolimits_{w\in R_q}\operatorname{SD}\bigl(
		\mathsf{Pre}(h,\mathcal B,w)&,
		D_{t_w+\Lambda_h,\sigma_r}\bigr)
		\leq\delta_{\mathsf{GPV}}+d\nu.
	\end{aligned}
	\]
	The comparison in the last line is before ring reduction; reduction
	modulo $q$ cannot increase the distance. For the first bound, every
	coset mass lies within the usual $(1\pm\epsilon)$ multiples of a
	common mass, which bounds its normalized image probability relative
	to uniform. For the second, multiplying the $d$ shifted-Gaussian
	normalizer ratios bounds the sampler's density ratio by
	$((1+\epsilon)/(1-\epsilon))^d$; the $d$ approximate scalar calls
	add at most $d\nu$ by coupling. These estimates are uniform in $w$,
	not merely averages over random targets.
	
	For a forward draw $r\leftarrow D_{\sigma_r}^{2}$, the conditional
	law given $f_h(r)=w$ is exactly that coset Gaussian. Replacing its
	image by uniform and then using $\mathsf{Pre}$ therefore costs at most
	\[
	\epsilon_{\mathsf{pre}}
	\leq\delta_{\mathsf{img}}+\delta_{\mathsf{GPV}}+d\nu
	\]
	per token, including its later opening. For $d=8192$ and
	$\epsilon=2^{-320}$, a scalar sampler meeting the \emph{design target}
	$\nu\leq2^{-340}$ would give $\epsilon_{\mathsf{pre}}<2^{-305}$.
	This is a conditional numerical budget, not a measurement of an
	implemented sampler. A faster Falcon sampler can replace the reference
	algorithm only with a guarantee for the same conditional distribution.
	
	\subsection{Composition over Keys and Tokens}
	Let $K\leq L-1$ count generated nondistinguished key records, and let
	$T_k$ count replaced tokens under record $k$, so $\sum_{k=1}^KT_k\leq T$.
	Let $\delta_{\mathsf{gen}}$ bound one implemented capped generator's
	distance from the ideal capped experiment, including its proposal
	sampling and numerical errors. Exact GS certification leaves no
	additional accepted-bad-basis event. Independent attempts give
	\[
	\begin{aligned}
		\epsilon_{\mathsf{badTD}}
		&\leq K\bigl((1-p_0)^{J_{\mathsf{gen}}}
		+\delta_{\mathsf{gen}}\bigr),\\
		\epsilon_{\mathsf{TD}}
		&\leq K\epsilon_{\mathsf{key}}
		+\sum_{k=1}^K T_k\epsilon_{\mathsf{pre}}
		+\epsilon_{\mathsf{badTD}}\\
		&\leq (L-1)\epsilon_{\mathsf{key}}
		+T\epsilon_{\mathsf{pre}}+\epsilon_{\mathsf{badTD}}.
	\end{aligned}
	\]
	The term $K\epsilon_{\mathsf{key}}$ accounts for the computational
	replacement of public labels.
	The term $\sum_{k=1}^{K}T_k\epsilon_{\mathsf{pre}}$ accounts for the
	statistical comparison of tokens and their later openings.
	The term $\epsilon_{\mathsf{badTD}}$ accounts for capped-generation
	failure and generator implementation error.
	These contributions are collected in $\epsilon_{\mathsf{TD}}$,
	with generation failure and implementation error counted only once. For $K\geq1$, choosing
	$J_{\mathsf{gen}}\geq(\kappa\ln2+\ln K)/p_0$ bounds the total ideal
	generation-failure probability by $2^{-\kappa}$. The per-image uniform
	bound permits adaptively chosen targets and repeated use of a key.
	The outer label-guessing factor $L$ in Theorem~1 still applies to
	$\epsilon_{\mathsf{TD}}$. Errors already included in
	$\delta_{\mathsf{gen}}$ or $d\nu$ are not charged again in the
	whole-run error for actual protocol samples.
	
	\subsection{Hardness Assumptions}
	The theorem is stated using this trapdoor contract, Hint-RLWE, and
	normalized ring-SIS. It does not require a separate reduction to a
	particular standard formulation of NTRU if screened-key
	indistinguishability is adopted explicitly as above. For this proposed
	NTRU-type generator, however, that is an additional NTRU-like
	decisional assumption, not something that can be omitted. Search-NTRU
	key-recovery costs do not alone establish it. Removing this dependency
	would require another generator satisfying the same contract, or a
	different simulation proof. The formal scheme's honest algorithms do
	not themselves invoke the generator, but the present proof needs it
	to simulate nondistinguished tokens and completed responses.
	
\end{document}
